% Заключение

В настоящей работе разработан и экспериментально обоснован метод адаптации принципа B-LoRA~\cite{frenkel2024blora} к архитектуре многомодального диффузионного трансформера FLUX.1~\cite{blackforest2024flux}. Предложена гипотеза о функциональной специализации трансформерных блоков модели: поздние блоки однопоточной обработки (блоки 30--37 из 38) преимущественно кодируют стилистические атрибуты изображения, тогда как ранние блоки двупоточной обработки (блоки 0--5 из 19) несут ответственность за семантическое содержание сцены. Для проверки гипотезы спроектированы три серии аблационных экспериментов: по диапазону обучаемых блоков~(A), рангу адаптера~(B) и числу шагов обучения~(C), а также выполнено сравнение с базовыми методами~--- Full-LoRA-FLUX, IP-Adapter-FLUX~\cite{ye2023ipadapter} и B-LoRA-SDXL~\cite{frenkel2024blora}~--- по метрикам CLIP-style, CLIP-content~\cite{radford2021clip}, FID и LPIPS.

\textbf{Основным результатом работы является подтверждение гипотезы о специализации блоков архитектуры FLUX.1}: конфигурация $\Theta_{\mathrm{style}}$ на блоках single-stream [30--37] при совместном обучении с $\Theta_{\mathrm{content}}$ на блоках double-stream [0--5] обеспечивает наилучшее соотношение показателей CLIP-style и CLIP-content относительно сравниваемых методов. Расширение диапазона стилевых блоков до [24--37] и [19--37] приводит к росту смешения стиля и содержания, свидетельствуя о неоднородности функциональной специализации в пространстве однопоточных блоков. Аблации по рангу показывают, что значение $r = 16$ является оптимальным: при меньшем ранге ($r = 4$) адаптер недостаточно захватывает стилистические детали, при большем ($r = 32$) наблюдается переобучение при малом обучающем наборе. Оптимальное число шагов дообучения составляет 500: меньшее число шагов не позволяет адаптеру в полной мере усвоить стиль, тогда как 1\,000 шагов сопровождаются деградацией контентной составляющей генерации.

\textbf{Практическая значимость} разработанного метода состоит в том, что он позволяет выполнять однофотографический перенос художественного стиля на базе модели FLUX.1 без модификации весов базовой модели и без сбора обширных обучающих выборок. Для обучения достаточно единственного изображения-референса; при инференсе подключается исключительно стилевой адаптер~$\Theta_{\mathrm{style}}$, а содержание генерируемой сцены произвольно задаётся текстовым промптом. Это обеспечивает сочетание высокой стилистической верности и гибкого контентного управления, практически значимое в задачах цифрового искусства, персонализированной генерации и промышленного дизайна.

Вместе с тем предложенный подход имеет ряд \textbf{ограничений}. Во-первых, разделение стиля и содержания остаётся неполным: при стилях с сильной пространственной структурой (например, пуантилизм или геометрический орнамент) стилевой адаптер частично захватывает контентные признаки, что снижает управляемость генерации. Во-вторых, обучение на одном изображении чувствительно к его репрезентативности: нетипичные для художника образцы могут приводить к нестабильному переносу и артефактам при промптах, далёких от содержания референса.

\textbf{Направления будущих исследований} включают несколько перспективных линий. Первое~--- систематическое исследование иных диапазонов блоков и их перекрывающихся комбинаций, включая промежуточные однопоточные блоки [19--29], которые в рамках настоящей работы не рассматривались как основные кандидаты для $\Theta_{\mathrm{style}}$. Второе~--- совместный инференс адаптеров $\Theta_{\mathrm{style}} + \Theta_{\mathrm{content}}$ с независимо подобранными источниками стиля и содержания, что позволит реализовать полноценную задачу разделённого style-content transfer, аналогично исходному методу B-LoRA для SDXL. Третье~--- перенос предложенного принципа поблочной специализации LoRA на другие архитектуры класса DiT и MM-DiT: в частности, на Stable Diffusion~3~\cite{esser2024sd3}, поскольку её бэкбон также построен на MM-DiT и может обнаруживать схожие закономерности специализации, что открывает возможность унифицированного межархитектурного подхода к разделению стиля и содержания в современных диффузионных трансформерах.
